\begin{frame}{채점 기준 (rubric)}
  \small
  \begin{tabular}{@{}p{3.2cm}p{1.2cm}p{9.9cm}@{}}
    \toprule
    항목 & 비중 & 좋은 답 \\
    \midrule
    문제 정의와 데이터 & 15\% & 행·라벨·지표·\textbf{두 층 베이스라인}이
      명확하고, \textbf{정형/비정형 판정}이 데이터 특성에 근거함 \\
    방법론 적절성 & 25\% & \textbf{GBDT와 비교해서} ``왜 이 딥러닝
      기법''의 근거가 Ch07·Ch10의 개념(차원·구조·공유)으로 설명됨 \\
    검증의 정직성 & 25\% & 3분할·train-fit·\textbf{test 한 번}·
      베이스라인·\textbf{에포크 학습 곡선}·val 튜닝이 보고서에서
      추적 가능 \\
    수식적 정당화 & 20\% & Ch09~15 공식(10.1 파라미터, 9.3 정규화,
      15 ELBO 중 $\geq$1개)으로 ``왜''를 설명했고, \textbf{관측(숫자)
      과 일치} \\
    결과의 정직성과 반성 & 15\% & GBDT보다 못 한 부분·학습이 안 된
      부분을 \textbf{숨기지 않고} Ch09~15의 개념으로 진단 \\
    \bottomrule
  \end{tabular}
  \vskip0.5em
  \begin{block}{8.1과의 차이}
    ``방법론 적절성''이 ``이 데이터에 맞는 \textit{고전 ML 모델}인가''
   에서 ``\textbf{딥러닝이 필요한가, 그리고 그 중 왜 \textit{이}
    기법인가}''로 확장 — 성공의 기준이 ``신경망을 돌렸는가''가 아니라
    ``필요했는지를 \textbf{GBDT와 비교해 판정했는가}''이므로 25\%.
    성능이 GBDT보다 0.5\%p 낮아도 ``64차원 정형이므로 GBDT가 상한
   이고, EMNIST(784차원)로 확대해 검증했다''는 논리까지 세우면
    높은 점수.
  \end{block}
\end{frame}
